\documentclass[10pt, aspectratio=169]{beamer}

% Configuração de Tema e Cores
\usetheme{Madrid}
\usecolortheme{default}
\setbeamertemplate{navigation symbols}{}

% Pacotes essenciais
\usepackage[utf8]{inputenc}
\usepackage[T1]{fontenc}
\usepackage[brazil]{babel}
\usepackage{amsmath,amsthm,amssymb}
\usepackage{graphicx}
\usepackage{listings}
\usepackage{xcolor}

% Configuração de estilo de código (Listings para JavaScript)
\definecolor{codegreen}{rgb}{0,0.6,0}
\definecolor{codegray}{rgb}{0.5,0.5,0.5}
\definecolor{codepurple}{rgb}{0.58,0,0.82}
\definecolor{backcolour}{rgb}{0.95,0.95,0.92}

\lstdefinestyle{mystyle}{
    backgroundcolor=\color{backcolour},   
    commentstyle=\color{codegreen},
    keywordstyle=\color{magenta},
    numberstyle=\tiny\color{codegray},
    stringstyle=\color{codepurple},
    basicstyle=\ttfamily\footnotesize,
    breakatwhitespace=false,         
    breaklines=true,                 
    captionpos=b,                    
    keepspaces=true,                 
    numbers=left,                    
    numbersep=5pt,                  
    showspaces=false,                
    showstringspaces=false,
    showtabs=false,                  
    tabsize=2
}
\lstset{style=mystyle}

% Informações da Capa
\title[Simulação N-Corpos e Eficiência Gráfica]{Simulação de Gravidade 3D: Da Renderização ao Desempenho}
\subtitle{Otimização Visual e Física com Particionamento Espacial (Octrees)}
\author[Herivelton Siqueira]{Herivelton Guilherme Alves de Siqueira}
\institute[IMPA Tech]{
    Bacharelado em Matemática da Tecnologia e Inovação\\
    IMPA Tech
}
\date{Junho de 2026}

\begin{document}

% Slide 1: Capa
\begin{frame}
    \titlepage
\end{frame}

% Slide 2: Introdução e Motivação
\begin{frame}{Objetivo: A Construção de um Mundo 3D do Zero}
    \begin{itemize}
        \item \textbf{A Proposta:} Desenvolver um motor físico e gráfico nativo, executado puramente via navegador (\textit{Canvas 2D}), capaz de simular e renderizar a dinâmica de corpos celestes em 3D.
        \item \textbf{O Desafio Visual:} Projetar coordenadas tridimensionais exatas do espaço "mundo" $(x, y, z)$ para o espaço bidimensional da tela do monitor $(x, y)$.
        \item \textbf{O Desafio de Desempenho (O Gargalo):} Garantir uma taxa de atualização estável de 60 quadros por segundo (FPS). Renderizar múltiplos corpos interagindo de forma verossímil exige extrema eficiência algorítmica.
    \end{itemize}
\end{frame}

% Slide 3: O Pipeline Gráfico e a Câmera
\begin{frame}{O Núcleo Visual: Projeção Perspectiva}
    \begin{block}{Simulando a Câmera Virtual}
        Na ausência de um motor comercial (como a Unity), o projeto implementa as etapas do \textit{Pipeline Gráfico Clássico} através de transformações com matrizes matemáticas.
    \end{block}
    
    \vspace{0.3cm}
    
    \begin{itemize}
        \item A transformação de perspectiva obedece à regra de que \textbf{objetos mais distantes devem convergir e ser renderizados em escalas menores}.
        \item Multiplicamos as coordenadas reais pelas Matrizes de Câmera, gerando a imersão de profundidade (3D).
        \item \textbf{Consequência Computacional:} O processamento e a renderização dessas posições e forças a cada \textit{frame} sobrecarrega a CPU. Precisamos de estratégias de otimização em CG para a cena escalar.
    \end{itemize}
\end{frame}

% Slide 4: Implementação da Projeção
\begin{frame}[fragile]{Implementação: Pipeline de Projeção (\textit{camera.js})}
    Como as matrizes transformam as coordenadas (Espaço $Mundo \rightarrow NDC$):
\begin{lstlisting}[language=Java]
// Transformacao linear combinada global (ViewProj)
window.vec4.transformMat4(clip, v4, this.viewProjMatrix);
const w = clip[3]; // Fator de escala da profundidade

if (w <= 0) return null; // Clipping Plane Frontal

// Divisao Perspectiva: Normalizacao para NDC [-1, 1]
const ndcX = clip[0] / w;
const ndcY = clip[1] / w;
const ndcZ = clip[2] / w;

// Mapeamento de Viewport (Conversao para pixels reais)
const sx = (ndcX * 0.5 + 0.5) * this.canvas.width;
const sy = (1.0 - (ndcY * 0.5 + 0.5)) * this.canvas.height;
\end{lstlisting}
    \textit{Note a Divisão Perspectiva ($/w$), crucial na Computação Gráfica para criar a ilusão de profundidade escalando pontos mais distantes.}
\end{frame}

% Slide 5: O Gargalo Computacional
\begin{frame}{O Problema da Força Bruta: Complexidade $\mathcal{O}(N^2)$}
    \begin{itemize}
        \item Em uma cena física interativa, calcular a gravidade e detectar colisões segue, por padrão, um modelo ingênuo de "todos-contra-todos".
        \item Para $N$ corpos celestes na tela, o computador realiza na ordem de $N^2$ cálculos por frame.
        \item Se dobrarmos a população de objetos na simulação, o custo de processamento \textbf{quadruplica}. Esse modelo quadrático inviabiliza testes reais em Computação Gráfica.
    \end{itemize}
    
    \begin{block}{A Necessidade Teórica}
        Para resolver isso de forma elegante, a teoria da Computação Gráfica dita o uso de ferramentas de \textbf{Estruturas de Dados de Particionamento Espacial}.
    \end{block}
\end{frame}

% Slide 6: Octree - A Solução de Otimização
\begin{frame}{A Estrutura Espacial: A Octree}
    \begin{itemize}
        \item \textbf{O Conceito Teórico:} Uma Octree é uma estrutura em árvore que subdivide recursivamente o volume 3D total do "mundo" em 8 sub-regiões (octantes) contínuas, garantindo que cada pequena área armazene apenas uma capacidade máxima de corpos (ex: 4 corpos por folha).
    \end{itemize}
    
    \begin{block}{Eficiência na Prática (Implementação no Projeto)}
        Como evidenciado através da malha tridimensional dinâmica e brilhante no visualizador (o esqueleto da Octree), a cena 3D é fatiada dinamicamente conforme os corpos se movem, atuando como um poderoso índice espacial.
    \end{block}

    \begin{itemize}
        \item \textbf{Resultado:} Redução abrupta no espaço de busca e varredura geométrica da cena.
    \end{itemize}
\end{frame}

% Slide 7: Implementação do Sub-Espaço
\begin{frame}[fragile]{Implementação: Subdivisão da Octree (\textit{octree.js})}
    O algoritmo cria recursivamente 8 filhos limitando os objetos por região:
\begin{lstlisting}[language=Java]
subdivide() {
    const nextSize = this.size / 2;
    const quarter = this.size / 4;
    // Cria 8 octantes combinando os eixos +/- X, Y, Z
    const signs = [-1, 1];
    for (const xSign of signs) {
        for (const ySign of signs) {
            for (const zSign of signs) {
                this.children.push(new OctreeNode(
                    this.cx + xSign * quarter,
                    this.cy + ySign * quarter,
                    this.cz + zSign * quarter,
                    nextSize, this.depth + 1, this.maxDepth
                ));
            } } }
\end{lstlisting}
\end{frame}

% Slide 8: Otimização Numérica: Algoritmo de Barnes-Hut
\begin{frame}{Aceleração em Larga Escala: Barnes-Hut e a Octree}
    \begin{block}{Critério de Aceitação Multipolar (MAC)}
        Ao usar a Octree, o algoritmo de Barnes-Hut avalia a seguinte relação geométrica: $\frac{s}{d} < \theta$.
        Onde $s$ é a largura da região, $d$ é a distância da região ao corpo focado, e $\theta$ é um limite fixo de precisão ($\approx 0.5$).
    \end{block}
    
    \vspace{0.3cm}
    
    \begin{itemize}
        \item \textbf{A Otimização Lógica:} Se uma aglomeração de planetas vizinhos satisfaz o MAC, eles estão distantes o bastante visual e fisicamente.
        \item Em vez de somar vetores par a par, o motor trata aquela área agrupada da Octree como um único \textbf{Centro de Massa Virtual}.
        \item \textbf{Eficiência Extrema:} A complexidade temporal decai de $\mathcal{O}(N^2)$ para assintóticamente $\mathcal{O}(N \log N)$. Isso garante fluidez no navegador nativo.
    \end{itemize}
\end{frame}

% Slide 9: Implementação do MAC
\begin{frame}[fragile]{Implementação: Avaliação MAC (\textit{octree.js})}
    Decisão algorítmica de interagir com o nó inteiro ou adentrar em seus filhos:
\begin{lstlisting}[language=Java]
// MAC: s/d < theta  ->  no suficientemente distante
if (dist > 0.001 && (this.size / dist) < theta) {
    // Trata todo o octante como um unico corpo virtual 
    // ancorado no centro de massa
    let accMag = (G_SCALED * this.totalMass) / distSq;
    
    const ax = accMag * (dx / dist); // X
    const ay = accMag * (dy / dist); // Y
    const az = accMag * (dz / dist); // Z
    
    body.addAcceleration(ax, ay, az);
} else {
    // No muito proximo -> desce na arvore (maior precisao)
    for (const child of this.children) 
        child.accumulateForce(body, G, theta, softening);
}
\end{lstlisting}
\end{frame}

% Slide 10: Colisões (Broad-Phase)
\begin{frame}{Eficiência Aplicada à Intersecção 3D (Colisões)}
    \begin{itemize}
        \item O teste de sobreposição de malhas e esferas é outro gargalo famoso nos Motores Gráficos 3D (\textit{Game Engines}).
    \end{itemize}
    
    \begin{block}{A Abordagem \textit{Broad-Phase} via Octree}
        Graças à Octree já calculada no ciclo do \textit{frame}, o algoritmo não permite o teste inútil de esferas que se encontram em cantos opostos da cena.
        Ele restringe estritamente a detecção mútua aos corpos que habitam o mesmo nó-folha da árvore no espaço delimitado da malha.
    \end{block}

    \begin{itemize}
        \item \textbf{Conclusão Visual (Narrow-Phase):} Detectada a intersecção euclidiana, opera-se a correção de posição (separação dos vértices para evitar o artefato gráfico dos corpos "engolindo" uns aos outros) e a devolução realista do rebote.
    \end{itemize}
\end{frame}

% Slide 11: Conclusões
\begin{frame}{Considerações Finais}
    \begin{itemize}
        \item O desenvolvimento deste projeto ilustra que \textbf{a excelência do desempenho (eficiência) e estruturação algorítmica é tão essencial quanto a teoria visual} em Computação Gráfica Interativa.
        \item A escolha estratégica da \textbf{Octree e Barnes-Hut} alinhada à matemática matricial nativa, comprova a eficácia de contornar gargalos quadráticos para atingir o \textit{Real-Time}.
        \item A implementação evidenciou, em código funcional e malha renderizável, como a ciência de particionamento e hierarquia espacial ditam o sucesso das renderizações modernas em ambientes tridimensionais escaláveis.
    \end{itemize}
    
    \vspace{0.5cm}
    \centering
    \textbf{Obrigado!}
\end{frame}

% Slide 12: Referências
\begin{frame}{Referências Bibliográficas}
    \begin{thebibliography}{9}
        \footnotesize
        
        \bibitem{gomes}
        GOMES, J.; VELHO, L. \textit{Fundamentos da Computação Gráfica}. IMPA, 2015.
        
        \bibitem{shirley}
        MARSCHNER, S.; SHIRLEY, P. \textit{Fundamentals of Computer Graphics}. 5. ed. AK Peters, 2021.
        
        \bibitem{samet}
        SAMET, H. \textit{Foundations of Multidimensional and Metric Data Structures}. Morgan Kaufmann, 2006.
        
        \bibitem{ericson}
        ERICSON, C. \textit{Real-Time Collision Detection}. Morgan Kaufmann, 2004.
        
    \end{thebibliography}
\end{frame}

\end{document}